Probe selection precedes primary scanning and selects inverted lists,
not database-vector survivors. For IVF centroid $\mu_i$, the selector
orders the FP32-valued computed routing scores
\begin{equation}
 s_i=\|\mu_i\|_2^2-2\langle q',\mu_i\rangle,
 \qquad (s_i,i)\ \text{ordered lexicographically}.
 \label{eq:5-probe-score}
\end{equation}
Omitting the common query norm preserves centroid ordering. In contrast
to primary squared-distance tables, these scores can be negative. The
required output is the first $n_{\rm sel}=\min(n_{\rm probe},n_{\rm list})$
centroid IDs in score order, using the smaller list ID at a tie, plus
the prefix offsets of their list lengths.

The dot products inherit the cuBLAS math mode: in the measured v69,
OpenAI rotation and centroid GEMMs use TF32 inputs with FP32 accumulation.
The selector's exactness concerns these computed scores. Its subsequent
FP32 \texttt{fma} and deterministic key ordering do not recover rounding
already introduced by the GEMM (\cref{sec:7-precision}).

\subsection{Radix selection followed by deterministic ordering}

The earlier bitonic selector repeatedly bitonic-sorted a growing candidate buffer.
The radix replacement keeps one block per query but finds the cutoff by
four most-significant-byte-first radix passes. Each pass counts only
keys matching the prefix selected so far, then locates the bin containing
the remaining order statistic. After four passes it has the cutoff key
and the number of equal-cutoff entries still required.

The float mapping canonicalises both signed zeros to positive zero,
then transforms their 32-bit representation $u$ by
\begin{equation}
 F(u)=\begin{cases}
       \mathop{\sim}u,&\text{if the sign bit is set},\\
       u\mathbin{|}2^{31},&\text{otherwise}.
      \end{cases}
 \label{eq:5-probe-key}
\end{equation}
For non-NaN floats this preserves numerical order, including negative
scores and infinities; canonicalisation makes equal zeros share a key.
This differs from the nonnegative-score encoding in
\cref{eq:5-packed-key}. NaNs are outside the selector's ordering contract.
Every pass recomputes scores with the same explicitly rounded
\texttt{fma} expression, so changing the pass does not change its keys.

The histogram has 256 bins and 32 lane replicas, requiring 32\,KiB for
its counters alone. Replicas separate same-warp atomic updates to a
popular bin; different warps can still contend. This exchanges shared
storage and reduction work for less concentrated histogram updating,
rather than eliminating all contention.

Collection first gathers every key below the cutoff. If too many keys
equal the cutoff, an ordered rescan uses ballots and prefix counts to
take exactly the required number in increasing list-ID order. Atomic
arrival order therefore does not choose the final tied subset.
For each of the $n_{\rm sel}$ selected entries, a rank-count step counts how many
selected keys precede it under $(s_i,i)$ and writes it to that rank.
This restores deterministic sorted order without sorting all centroids.

Finally, a block prefix scan produces offsets of the selected list
lengths and their total candidate count. When $n_{\rm probe}>n_{\rm list}$,
the implementation pads the ID array by repeating the final selected ID
but assigns those extra entries zero length. They do not scan the same
list again. Nonempty centroid sets and a positive requested probe count
are assumed.

\subsection{Costs and separate correctness contracts}

The histogram passes and collection revisit the centroid-score row;
rank counting costs $O(n_{\rm sel}^2)$ comparisons. The algorithm is consequently
not a uniformly linear-cost substitute for every probe count. Large
$n_{\rm sel}$ increases both ranking work and selected-entry shared storage,
while the histogram reserves space even for small centroid sets.
Radix selection is an established selection technique; the contribution
here is its deterministic integration and measured execution trade-off
for this query pipeline.



The primitive tests cover many ties, all-equal values, mixed signed
zeros, mostly infinite values, extreme magnitudes and requested probe
counts exceeding the list count. Their outcomes and real-query checks
are reported in \cref{sec:7-4}.


The optional FP16 routing changes score precision and is evaluated as
a separate approximation in \cref{sec:7-4}, alongside the cross-process
pipeline-identity checks.
